% ─────────────────────────────────────────────────────────────────────────────
% I-Cache Controller FSM — tikz state diagram and state descriptions
% Source: ip_icache/src/asICache.sv
% ─────────────────────────────────────────────────────────────────────────────

The complete Moore-FSM of \texttt{asICache} is shown in
Fig.~\ref{fig:asICacheFsm01}.
Outputs depend only on the current state (state nodes contain key active
signals); input conditions label the transitions.
All signal names correspond directly to the \texttt{asICache} port list
(Tables~\ref{tab:asICacheCtrlPorts}~and~\ref{tab:asICacheCtrlPorts2}).

\begin{figure}[H]
\centering
\begin{tikzpicture}[shorten >=1pt, node distance=3.2cm, on grid, auto,
  every state/.style={fill=blue!10}, initial text=]

  % ── State nodes ──────────────────────────────────────────────────────────

  \node[state with output, initial above, accepting] (idle)
    { \texttt{IDLE}
      \nodepart{lower}
      \begin{tiny}$\begin{array}{l}
        \mathit{stall}{=}0 \\
        \mathit{rvalid}{=}0
      \end{array}$\end{tiny}
    };

  \node[state with output] (flush) [right=of idle]
    { \texttt{FLUSH}
      \nodepart{lower}
      \begin{tiny}$\begin{array}{l}
        \mathit{stall}{=}1
      \end{array}$\end{tiny}
    };

  \node[state with output] (lookup) [below=of idle]
    { \texttt{LOOKUP}
      \nodepart{lower}
      \begin{tiny}$\begin{array}{l}
        \mathit{stall}{=}1
      \end{array}$\end{tiny}
    };

  \node[state with output] (miss) [below=of lookup]
    { \texttt{MISS}
      \nodepart{lower}
      \begin{tiny}$\begin{array}{l}
        \mathit{stall}{=}1 \\
        \mathit{arvalid}{=}1
      \end{array}$\end{tiny}
    };

  \node[state with output] (fill) [below=of miss]
    { \texttt{FILL}
      \nodepart{lower}
      \begin{tiny}$\begin{array}{l}
        \mathit{stall}{=}1 \\
        \mathit{rready}{=}1
      \end{array}$\end{tiny}
    };

  \node[state with output] (respond) [below=of fill]
    { \texttt{RESPOND}
      \nodepart{lower}
      \begin{tiny}$\begin{array}{l}
        \mathit{rvalid}{=}1 \\
        \mathit{stall}{=}0
      \end{array}$\end{tiny}
    };

  \node[state with output] (err) [right=of fill]
    { \texttt{ERR}
      \nodepart{lower}
      \begin{tiny}$\begin{array}{l}
        \mathit{err}{=}1 \\
        \mathit{stall}{=}0
      \end{array}$\end{tiny}
    };

  % ── Transitions ───────────────────────────────────────────────────────────

  \path[->]
    % IDLE arcs
    (idle)
      edge [in=100,out=130,loop]     node [left] {\scriptsize $\neg\mathit{req}{\wedge}\neg\mathit{flush}$} (idle)
      edge [bend right] node [left] {\scriptsize $\mathit{ic\_req}$}                                                         (lookup)
      edge [bend right] node [above] {\scriptsize $\mathit{ic\_flush}$}                                                 (flush)

    % LOOKUP arcs
    (lookup)
      edge [bend right=35] node [right] {\scriptsize $\mathit{hit}$}        (idle)
      edge                                 node [left] {\scriptsize $\neg\mathit{hit}$} (miss)

    % MISS arcs
    (miss)
      edge [loop left] node [left] {\scriptsize $\neg\mathit{arready}$} (miss)
      edge                     node [left] {\scriptsize $\mathit{arready}$}           (fill)

    % FILL arcs
    (fill)
      edge [loop left]     node [left] {\scriptsize $\mathit{rvalid}{\wedge}\neg\mathit{rlast}$}                                              (fill)
      edge                          node [left] {\scriptsize $\mathit{rvalid}{\wedge}\mathit{rlast}{\wedge}\mathit{rresp}{=}0$} (respond)
      edge [bend right] node [above] {\scriptsize $\mathit{rvalid}{\wedge}\mathit{rresp}{\neq}0$}                                  (err)

    % RESPOND arc back to IDLE (long, bend right to avoid crossing nodes)
    (respond)
      edge [bend left=90] node [right] {\scriptsize $\mathbf{1}$} (idle)

    % ERR arc back to IDLE
    (err)
      edge [bend right=30] node [above right] {\scriptsize $\mathbf{1}$} (idle)

    % FLUSH arcs
    (flush)
      edge [loop right]          node [right] {\scriptsize $\mathit{cnt}{<}31$}   (flush)
      edge [bend right=25] node [below] {\scriptsize $\mathit{cnt}{=}31$} (idle);

\end{tikzpicture}
\caption{Moore-FSM: \texttt{asICache} Instruction Cache Controller
  (\texttt{ip\_icache/src/asICache.sv}).
  \textbf{State nodes} show active outputs only; inactive signals are
  zero.
  $\mathit{cnt}$ is the 5-bit set counter \texttt{flush\_cnt\_r}
  (counts 0\ldots{}31 during cache-line invalidation).
  The self-loop label \textbf{1} on \textsc{Respond} and \textsc{Err}
  denotes an unconditional one-cycle transition.}
\label{fig:asICacheFsm01}
\end{figure}

\paragraph{State-by-state description of \texttt{asICache}:}

\begin{description}

  \item[\texttt{IDLE}]
    The cache is quiescent.
    \texttt{ic\_stall\_o} is de-asserted; the CPU may issue a new request.
    On each clock cycle one of three branches is taken:
    \begin{itemize}
      \item $\mathit{ic\_flush\_i}{=}1$: assert \texttt{ic\_stall\_o},
        initialise \texttt{flush\_cnt\_r} to zero, transition to
        \textsc{Flush}.
      \item $\mathit{ic\_req\_i}{=}1$: latch \texttt{ic\_addr\_i} into
        \texttt{lk\_tag\_r}, \texttt{lk\_idx\_r}, \texttt{lk\_instr\_r};
        present \texttt{req\_idx\_s} to the data and tag SRAM read ports
        (registered output will be valid in \textsc{Lookup});
        assert \texttt{ic\_stall\_o}; transition to \textsc{Lookup}.
      \item Otherwise: remain in \textsc{Idle}.
    \end{itemize}

  \item[\texttt{LOOKUP}]
    The SRAM registered output (\texttt{data\_rdata\_s},
    \texttt{tag\_rdata\_s}) is now valid.
    The combinatorial valid array \texttt{valid\_rdata\_s} is read via
    \texttt{lk\_idx\_r}.
    Hit detection computes \texttt{hit\_vec\_s[w] = valid[w] \& (tag[w]
    == lk\_tag\_r)} for all ways $w \in \{0,1,2,3\}$:
    \begin{itemize}
      \item \textbf{Hit} ($\mathit{hit\_s}{=}1$): extract the 32-bit
        instruction word from the hit way using
        \texttt{lk\_instr\_r[4:2]} as the selector; assert
        \texttt{ic\_rdata\_o} and \texttt{ic\_rvalid\_o}; update PLRU;
        de-assert \texttt{ic\_stall\_o}; transition to \textsc{Idle}.
        Total latency: \textbf{2 cycles} (IDLE + LOOKUP).
      \item \textbf{Miss} ($\mathit{hit\_s}{=}0$): select victim way via
        PLRU (prefer invalid way); latch \texttt{fill\_way\_r};
        drive the AXI4 AR channel with the cache-line-aligned miss
        address (\texttt{lk\_line\_addr\_r}); assert
        \texttt{ar\_valid\_r}; transition to \textsc{Miss}.
    \end{itemize}

  \item[\texttt{MISS}]
    The AXI4 AR channel handshake is in progress.
    \texttt{arvalid} is held high; the FSM waits until the memory
    arbiter (\texttt{asMemArb}) asserts \texttt{arready}.
    On acknowledgement: clear \texttt{ar\_valid\_r}; transition to
    \textsc{Fill}.
    Stall remains asserted throughout.

  \item[\texttt{FILL}]
    AXI4 R-channel beats are received.
    \texttt{rready} is held high (combinatorial from state).
    For each beat with $\mathit{rvalid}{=}1$:
    \begin{itemize}
      \item If $\mathit{rresp}{\neq}0$: assert \texttt{ic\_err\_o};
        de-assert stall; transition to \textsc{Err}.
      \item If $\mathit{rlast}{=}0$: write the beat into
        \texttt{fill\_buf\_r} at position \texttt{beat\_r}; increment
        \texttt{beat\_r}; remain in \textsc{Fill}.
      \item If $\mathit{rlast}{=}1$ and $\mathit{rresp}{=}0$:
        write the full line (data + tag + valid) to the SRAM;
        update PLRU; latch the requested instruction word into
        \texttt{rdata\_r}; transition to \textsc{Respond}.
    \end{itemize}

  \item[\texttt{RESPOND}]
    One-cycle response state.
    \texttt{ic\_rvalid\_o} is pulsed high; \texttt{ic\_stall\_o} is
    de-asserted; the FSM returns to \textsc{Idle}.
    This extra cycle exists because the SRAM write (in \textsc{Fill})
    and the read-data extraction occur in the same cycle, and the
    CPU must sample valid data on the following edge.

  \item[\texttt{ERR}]
    One-cycle error state.
    \texttt{ic\_err\_o} is pulsed high; \texttt{ic\_stall\_o} is
    de-asserted; the FSM returns to \textsc{Idle}.
    The filled cache line is \emph{not} written (valid bit stays clear).

  \item[\texttt{FLUSH}]
    Cache invalidation loop.
    In each cycle the combinatorial write-port decode
    (\texttt{valid\_flush\_en\_s}) clears all four valid bits for set
    index \texttt{flush\_cnt\_r}, and PLRU bits for that set are zeroed.
    After 32 cycles (all 32 sets), \texttt{ic\_flush\_done\_o} is
    pulsed and the FSM returns to \textsc{Idle} with stall de-asserted.

\end{description}
